%% ============================================================
%% SECTION 2: BACKGROUND — RS DOMAIN AND MLLM FOUNDATIONS
%% EarthVision-LM Survey — Artificial Intelligence Review
%% Template: sn-jnl (Springer Nature) | sn-mathphys-num
%% Rules: booktabs, TikZ, numbered [N] cites,
%%        every float cited before appearance, NO fabricated data
%% ============================================================

\section{Background}\label{sec2}

This section establishes the domain-specific context necessary to understand
why RS-MLLMs require architectures, datasets, and evaluation protocols
substantially different from those of general-purpose MLLMs. We first survey
the characteristics of RS imagery that make it a uniquely challenging domain
for language-grounded vision models, then review the MLLM foundations on which
RS-MLLMs are built, and conclude with the systematic literature review
methodology used in this survey.

%% ──────────────────────────────────────────────────────────────
\subsection{Remote Sensing Imagery: Domain Characteristics}\label{sec2:rs}
%% ──────────────────────────────────────────────────────────────

Remote sensing imagery is acquired by sensors mounted on satellites, airborne
platforms, or unmanned aerial vehicles. The resulting images differ from
ground-level photographs along every dimension that vision models exploit:
viewpoint, scale, spectral content, object orientation, and scene density.
Table~\ref{tab:rs_challenges} summarises the six domain-specific challenges
identified in this survey and maps each to the RS-MLLM design dimension it
motivates.

\begin{table}[h]
\caption{Six RS domain challenges and their downstream consequences for
RS-MLLM design. Each challenge motivates one dimension of the five-dimensional
design space introduced in Sect.~\ref{sec3}.}\label{tab:rs_challenges}
\small\setlength\tabcolsep{4pt}
\begin{tabular*}{\textwidth}{@{\extracolsep\fill}p{2.8cm}p{4.2cm}p{3.8cm}}
\toprule
\textbf{Challenge} & \textbf{Description} & \textbf{Design-Space Impact} \\
\midrule
Nadir-view domain gap &
  Ground-level visual priors (CLIP, BLIP-2) fail for top-down imagery &
  Dim~$\mathcal{E}$: encoder adaptation \\[3pt]
Extreme resolution range &
  VHR imagery (0.3--1\,m/px) produces token sequences exceeding LLM context &
  Dim~$\mathcal{C}$: token-budget connectors \\[3pt]
Object orientation ambiguity &
  RS objects appear at arbitrary rotation angles; HBB is a poor fit &
  Dim~$\mathcal{S}$: OBB spatial output \\[3pt]
Sensor heterogeneity &
  Optical, SAR, HSI, LiDAR encode physically incompatible modalities &
  Dim~$\mathcal{M}$: modality scope \\[3pt]
Compute constraints &
  Academic RS groups lack the GPU budget for full LLM fine-tuning &
  Dim~$\mathcal{L}$: freezing strategy \\[3pt]
Annotation scarcity &
  Expert RS annotation is costly; long-tail categories are under-represented &
  Training data and PEFT strategy \\
\botrule
\end{tabular*}
\end{table}

Figure~\ref{fig:rs_challenges} illustrates the six challenges visually,
connecting each to the RS imagery properties that produce it.

\begin{figure}[h]
\centering
\begin{tikzpicture}[font=\small, >=Stealth,
    ch/.style={draw, rounded corners=3pt, fill=#1!16,
               minimum width=3.8cm, minimum height=0.72cm,
               align=center, font=\small},
    arr/.style={->, thick, gray!55, shorten >=2pt, shorten <=2pt}]

  %% Central RS image icon
  \fill[gray!22] (0,-0.4) rectangle (2.0,1.6);
  \fill[gray!35] (0.1,-0.3) rectangle (1.9,1.5);
  %% Simulated RS scene features
  \fill[teal!30] (0.2,0.2) rectangle (0.8,0.8);   % field
  \fill[gray!55] (1.0,0.6) rectangle (1.7,1.3);   % building
  \fill[blue!25] (0.1,-0.2) rectangle (0.6,0.15);  % water
  %% Rotated object (aircraft)
  \fill[gray!65, rotate around={40:(1.3,0.1)}]
    (1.0,0.0) rectangle (1.6,0.2);
  \node[font=\scriptsize\bfseries, gray!80] at (1.0,2.0)
    {RS Nadir Image};

  %% Six challenge boxes — radial layout
  %% Top
  \node[ch=blue]   (c1) at (1.0, 3.6)
    {C1: Nadir-View\\Domain Gap};
  %% Top-right
  \node[ch=orange] (c2) at (5.2, 2.5)
    {C2: Extreme\\Resolution Range};
  %% Bottom-right
  \node[ch=red]    (c3) at (5.2,-1.0)
    {C3: Object\\Orientation};
  %% Bottom
  \node[ch=purple] (c4) at (1.0,-2.5)
    {C4: Sensor\\Heterogeneity};
  %% Bottom-left
  \node[ch=teal]   (c5) at (-3.2,-1.0)
    {C5: Compute\\Constraints};
  %% Top-left
  \node[ch=green!60!black] (c6) at (-3.2, 2.5)
    {C6: Annotation\\Scarcity};

  %% Arrows from central image to each challenge
  \draw[arr] (1.0,2.0) -- (c1.south);
  \draw[arr] (2.0,0.8) -- (c2.west);
  \draw[arr] (2.0,0.1) -- (c3.west);
  \draw[arr] (1.0,-0.4)-- (c4.north);
  \draw[arr] (0.0,0.1) -- (c5.east);
  \draw[arr] (0.0,0.8) -- (c6.east);

  %% Design-space labels (small, below each box)
  \node[blue!70, font=\tiny\bfseries] at (1.0,2.80)
    {$\rightarrow$ Dim $\mathcal{E}$: Encoder};
  \node[orange!80,font=\tiny\bfseries] at (5.2,1.75)
    {$\rightarrow$ Dim $\mathcal{C}$: Connector};
  \node[red!70,   font=\tiny\bfseries] at (5.2,-1.75)
    {$\rightarrow$ Dim $\mathcal{S}$: Spatial};
  \node[purple!70,font=\tiny\bfseries] at (1.0,-3.28)
    {$\rightarrow$ Dim $\mathcal{M}$: Modality};
  \node[teal!80,  font=\tiny\bfseries] at (-3.2,-1.75)
    {$\rightarrow$ Dim $\mathcal{L}$: LLM};
  \node[green!60!black,font=\tiny\bfseries] at (-3.2,1.72)
    {$\rightarrow$ Training Data};
\end{tikzpicture}
\caption{Six RS domain challenges (C1--C6) and their mapping to RS-MLLM
design dimensions. Each challenge arises directly from the physical
properties of RS imagery acquisition (nadir viewpoint, sensor physics,
annotation economics) and motivates a specific architectural or training
design choice analysed in Sect.~\ref{sec3}.}\label{fig:rs_challenges}
\end{figure}

\subsubsection{C1: Nadir-View Domain Gap}\label{sec2:nadirview}

All mainstream vision encoders used in MLLMs---CLIP~\cite{clip2021},
EVA-CLIP~\cite{fang2023evaclip}, DINOv2~\cite{oquab2023dinov2}---were
pre-trained on internet images depicting objects from human eye-level
perspectives. RS nadir imagery presents the same physical objects from
directly above, producing visual appearances that are qualitatively different
from the training distribution. A passenger aircraft viewed from 500\,m altitude
appears as a cross-shaped grey silhouette against a tarmac background; from
ground level, the same aircraft presents fuselage, wings, windows, and livery
colours that human photographers routinely capture. The cross-shaped aerial
appearance has no representation in CLIP's training data, creating a fundamental
domain gap~\cite{long2021creating}. This manifests as the \emph{blind-pair
problem}: two semantically distinct RS tiles---for example, a dense urban
neighbourhood and an industrial complex---can receive near-identical CLIP
embeddings because neither matches any dominant training concept~\cite{bommasani2021}.

The nadir-view gap also invalidates standard RS object detection datasets
when used for MLLM fine-tuning. Descriptions generated from CLIP embeddings
for RS images often misidentify land-cover categories, producing noisy
instruction data that propagates errors into the trained
RS-MLLM~\cite{ddfav2025}.

\subsubsection{C2: Extreme Resolution Range}\label{sec2:resolution}

RS sensors span an enormous resolution range: spaceborne optical sensors such
as Sentinel-2 operate at 10--60\,m ground sampling distance (GSD) per pixel;
commercial VHR satellites (WorldView-3, Pl{\'e}iades) reach 0.3\,m GSD; airborne
and UAV sensors achieve sub-centimetre GSD. At the VHR end, a single tile
covering $2\!\times\!2$\,km at 0.3\,m GSD contains approximately
$6700\!\times\!6700$ pixels---three orders of magnitude more pixels than the
$224\!\times\!224$ input size of standard CLIP encoders.

Naive rescaling discards the fine spatial detail that distinguishes car models,
ship classes, or aircraft types. Tiling at native resolution preserves detail
but produces thousands of patches per image, overwhelming LLM context windows.
The MLLM connector must therefore implement token-efficient strategies (perceiver
resampler, token compression) specifically designed for the RS resolution
challenge~\cite{uhrbat2025, geollava2025}.

\subsubsection{C3: Object Orientation Ambiguity}\label{sec2:orientation}

Objects captured from nadir appear at arbitrary rotation angles determined by
their heading at acquisition time, not by any camera pose that a photographer
controls. A ship travelling north-northeast, an aircraft taxiing at 27\,degrees
heading, and a vehicle parked diagonally across two parking bays all produce
rotated object silhouettes in the RS image. Horizontal bounding boxes (HBB),
aligned to image axes, enclose substantial background area for any object
deviating from horizontal or vertical orientation. For a typical aircraft
rotation of 35--45\,degrees, the HBB waste factor reaches 55--70\,\% of the
box area~\cite{xia2018dota, ding2021object}.

Oriented bounding boxes (OBB), parameterised as $(x,\,y,\,w,\,h,\,\theta)$,
solve this problem by rotating the box to match the object heading, but require
a model capable of regressing the angular parameter $\theta$---a capability
absent from all general MLLM architectures and present in only one published
RS-MLLM (GeoChat~\cite{geochat2024}, as established in Sect.~\ref{sec3:spatial}).

\subsubsection{C4: Sensor Heterogeneity}\label{sec2:sensor}

Operational RS depends on multiple sensor modalities that each encode distinct
physical quantities. Optical sensors measure solar reflectance in discrete
spectral bands; the resulting images are the most visually interpretable but
are degraded by cloud cover, atmospheric scattering, and illumination variation.
SAR sensors transmit microwave pulses and measure the backscattered signal,
encoding surface roughness, dielectric constant, and double-bounce geometry
rather than colour~\cite{sentinel12014, schmitt2019fusion}. Hyperspectral sensors
resolve 200--400 contiguous narrow spectral bands enabling material identification
through spectral fingerprinting inaccessible to RGB or even multispectral
sensors~\cite{hong2021multimodal}.

Each modality has complementary strengths: SAR provides cloud-penetrating
all-weather coverage critical for disaster response; HSI enables crop-type
classification and mineral mapping impossible from optical data alone. A
geospatially intelligent RS-MLLM must ultimately reason across all three
modalities. However, most currently documented RS-MLLM vision encoders
inherit predominantly optical/RGB-oriented visual pre-training, while
modality-specific SAR and HSI representation learning remains substantially
underdeveloped, creating a sensor heterogeneity gap that limits the
applicability of RGB-trained features to SAR and HSI data~\cite{hmbench2026}.

\subsubsection{C5: Compute Constraints}\label{sec2:compute}

Large language model fine-tuning at the scale of GPT-4~\cite{openai2023gpt4}
or LLaMA-2~\cite{llama2023} requires thousands of GPU-hours on high-memory
accelerators. The RS academic community---primarily university research groups
rather than industry laboratories---typically operates under 2--8 A100 GPU
budgets. Full fine-tuning of even a 7B-parameter backbone is rarely feasible.
Parameter-efficient fine-tuning (PEFT) techniques such as
LoRA~\cite{hu2022lora}, prefix tuning, and adapter modules enable meaningful
RS domain adaptation within these constraints, at the cost of reduced model
flexibility.

This compute gap has a secondary effect: RS instruction datasets are
necessarily smaller than those used for general MLLMs. GeoChat-Instruct
contains 318K instruction pairs~\cite{geochat2024}; SkyEye-968K
contains approximately 968K pairs~\cite{skyeyegpt2025}; Earth-OneVision
trains on 34M pairs~\cite{earthonevision2026}---but general MLLM datasets
(LLaVA-1.5, ShareGPT4V) routinely exceed 100M examples. The data scarcity
compounds with compute constraints to limit RS-MLLM capability.

\subsubsection{C6: Annotation Scarcity and Long-Tail Distribution}\label{sec2:longtail}

Expert RS annotation requires spectral domain knowledge (to identify crop
varieties, mineral deposits, or structural damage types) that is expensive
and scarce. Unlike natural image datasets where crowdsourced annotation
(ImageNet, COCO) scales annotation effort, RS annotation requires
domain experts---remote sensing scientists, geographers, or military
image analysts~\cite{long2021creating}.

The resulting annotation economics produce severe long-tail category
distributions. In DOTA~\cite{xia2018dota}, the most populated category
(``large vehicle'') has 29,074 instances; the least populated (``helicopter'')
has 548---a 53-fold imbalance within a single dataset. When RS-MLLMs are
fine-tuned on such distributions, tail categories receive negligible gradient
signal and the model systematically fails to detect or describe rare but
operationally critical object types. Our prior work~\cite{odfe2025}
addresses this long-tail problem in the specialised OBB detection setting
through optimisation-driven feature engineering; extending analogous
strategies to the RS-MLLM training setting represents an open research
direction.

%% ──────────────────────────────────────────────────────────────
\subsection{MLLM Foundations}\label{sec2:mllm}
%% ──────────────────────────────────────────────────────────────

Multimodal large language models extend transformer-based LLMs~\cite{vaswani2017attention}
with the ability to process visual inputs. The architectural blueprint
established by BLIP-2~\cite{blip22023} and LLaVA~\cite{llava2023}---a frozen
visual encoder connected to a frozen or lightly fine-tuned LLM via a learned
vision-language connector---has been adopted by virtually every RS-MLLM in
the literature. Figure~\ref{fig:mllm_blueprint} illustrates this canonical
architecture.

\begin{figure}[h]
\centering
\begin{tikzpicture}[font=\small, >=Stealth,
    block/.style={draw, rounded corners=3pt, fill=#1,
                  minimum width=2.6cm, minimum height=0.85cm,
                  align=center},
    frozen/.style={draw, rounded corners=3pt, fill=#1,
                   minimum width=2.6cm, minimum height=0.85cm,
                   align=center, dashed}]

  %% Input
  \node[block=gray!15] (img) at (0,0) {RS Image\\$H\!\times\!W\!\times\!C$};
  \node[block=gray!15] (txt) at (0,-1.5) {Text Prompt\\(instruction)};

  %% Visual encoder
  \node[frozen=blue!15] (enc) at (3.2,0)
    {Vision Encoder\\(CLIP\,/\,ViT)\\{\scriptsize\textit{frozen / adapted}}};

  %% Connector
  \node[block=orange!20] (conn) at (6.2,0)
    {Connector\\(MLP\,/\,Q-Former\\/\,Perceiver)};

  %% LLM
  \node[frozen=green!15] (llm) at (9.0,-0.75)
    {LLM Backbone\\(Vicuna\,/\,LLaMA\\InternLM\,/\,Qwen)\\{\scriptsize\textit{frozen\,+\,LoRA}}};

  %% Output
  \node[block=teal!20] (out) at (9.0,1.2)
    {Output\\(Text\,/\,HBB\\OBB\,/\,Seg)};

  %% Arrows
  \draw[->,thick] (img.east)  -- (enc.west);
  \draw[->,thick] (enc.east)  -- (conn.west);
  \draw[->,thick] (conn.east) -- (llm.north west);
  \draw[->,thick] (txt.east)  -- (llm.south west);
  \draw[->,thick] (llm.north) -- (out.south);

  %% Token labels
  \node[gray!60, font=\scriptsize] at (1.6, 0.22) {image};
  \node[gray!60, font=\scriptsize] at (1.6,-0.15) {pixels};
  \node[orange!70,font=\scriptsize] at (4.7, 0.22) {visual};
  \node[orange!70,font=\scriptsize] at (4.7,-0.15) {tokens};
  \node[teal!70,  font=\scriptsize] at (7.6, 0.22) {fused};
  \node[teal!70,  font=\scriptsize] at (7.6,-0.15) {tokens};

  %% Freeze annotation
  \draw[blue!50, dashed, rounded corners=3pt]
    (1.85,-0.55) rectangle (4.55, 0.55);
  \node[blue!50, font=\tiny] at (3.2,-0.72) {typically frozen};
  \draw[green!50!black, dashed, rounded corners=3pt]
    (7.6,-1.70) rectangle (10.4, 0.35);
  \node[green!50!black, font=\tiny] at (9.0,-1.87) {frozen\,+\,LoRA or partial unfreeze};
\end{tikzpicture}
\caption{Canonical MLLM architecture adopted by virtually all RS-MLLMs.
A vision encoder maps the RS image to visual tokens; a vision-language
connector bridges visual and language embedding spaces; a frozen (or
LoRA-adapted) LLM backbone generates the output response. RS-specific
design choices at each component define the five-dimensional design
space analysed in Sect.~\ref{sec3}.}\label{fig:mllm_blueprint}
\end{figure}

\paragraph{Visual Encoder.}
CLIP~\cite{clip2021} and its variants (EVA-CLIP, RemoteCLIP~\cite{remoteclip2024},
GeoRSCLIP~\cite{georscclip2024}) encode images as sequences of patch tokens via
a ViT~\cite{dosovitskiy2021vit} architecture trained with contrastive
image--text objectives. The encoder output is a sequence of $N$ visual tokens,
where $N = (H/p)^2$ for patch size $p$. For RS images at VHR resolution,
$N$ can reach tens of thousands.

\paragraph{Vision-Language Connector.}
The connector maps visual tokens to the LLM token embedding space.
BLIP-2~\cite{blip22023,li2023blip} popularised the Q-Former bottleneck, which uses
a fixed set of learnable query tokens to cross-attend to visual features.
LLaVA~\cite{llava2023} instead uses a simple linear projection (MLP),
trading information preservation for simplicity. Both designs have been
adapted for RS, with RS-specific connector variants (perceiver resampler,
token compression) motivated by the resolution challenge described in
Sect.~\ref{sec2:resolution}.

\paragraph{LLM Backbone.}
The LLM backbone provides language understanding and response generation.
RS-MLLMs predominantly use open-weight LLMs (Vicuna~\cite{vicuna2023},
LLaMA-2~\cite{llama2023}, InternLM~\cite{internlm2023}, Qwen~\cite{qwen2023})
rather than proprietary models, enabling fine-tuning within academic compute
budgets. LoRA~\cite{hu2022lora} reduces trainable parameters by more than
95\,\% by inserting low-rank adapter matrices into the attention layers, making
RS domain adaptation feasible on 2--4 NVIDIA A100 GPUs.

%% ──────────────────────────────────────────────────────────────
\subsection{Related Work}\label{sec2:related}
%% ──────────────────────────────────────────────────────────────

We position EarthVision-LM against six categories of prior work, identifying
the specific analytical capability absent from each.

\textbf{(0) Broad geospatial LLM systematic review in AIR (February 2026).}
Dorabantu \& Badea~\cite{dorabantu2026geollm} published a systematic review in
\emph{this journal} (Artificial Intelligence Review, February 2026) covering geospatial
reasoning and awareness in large language models broadly---including GeoLLMs,
GIS integration, spatial reasoning, and multimodal geospatial systems.
\emph{Gap}: Covers geospatial LLMs in general; does not focus on RS-specific architectures
(encoder adaptation, connector design, LLM strategy for RS); no sensor-physics analysis;
no SAR/HSI modality gap analysis; no hallucination taxonomy; no RS-MLLM-specific PRISMA review.
\emph{Relationship}: The present survey is the RS-architecture-specific complement to
Dorabantu \& Badea's broader GeoLLM perspective. Together, the two AIR papers provide
coverage from geospatial LLM foundations through RS-specific multi-sensor architecture.

\textbf{(1) RS VLM task-level reviews.}
Li~et~al.~\cite{li2024survey} survey RS vision-language models organised by
task (captioning, VQA, grounding) rather than by architectural design.
\emph{Gap}: No formal architectural design-space framework; no hallucination
taxonomy; no modality gap analysis; no PRISMA methodology.

\textbf{(2) RS-MLLM training and benchmark surveys.}
Weng~et~al.~\cite{weng2025isprs} survey RS-MLLMs in ISPRS Journal, focusing
on training datasets and benchmark evaluation.
\emph{Gap}: No architectural taxonomy; no sensor-physics analysis; no
hallucination taxonomy; no systematic modality coverage analysis.

\textbf{(3) Concurrent RS-MLLM diagnostic survey (July 2026).}
Ma~et~al.~\cite{ma2026dsorsp} (arXiv:2607.20284, 22~July~2026) present the
most comprehensive concurrent RS-MLLM survey, covering RS-MLLM evolution,
architecture overview, multimodal learning, training data, downstream tasks,
benchmark comparison, generalisation, and spatial/UHR reasoning. Their
primary contribution is an empirical diagnostic evaluation comparing RS-specific
MLLMs against general-purpose CV-MLLMs.
\emph{Gap}: No formal 5-dimensional architectural design-space taxonomy; no
sensor-physics analysis of SAR/HSI barriers; no multi-type hallucination
taxonomy; no PRISMA-compliant search protocol.
\emph{Relationship}: Explicitly complementary---Ma~et~al.\ establish
\emph{what} the gaps are empirically; EarthVision-LM explains \emph{why}
they exist architecturally.

\textbf{(4) General MLLM surveys.}
General surveys of MLLMs~\cite{yin2024survey,zhang2024vision} cover the
broader MLLM landscape including LLaVA, InstructBLIP, and Flamingo families.
\emph{Gap}: No RS-specific content; no nadir-view domain analysis; no OBB
spatial output; no SAR or HSI modality coverage; not applicable to the
RS operational context.

\textbf{(5) Broad deep learning in RS surveys.}
Zhu~et~al.~\cite{zhu2019sn}, Ma~et~al.~\cite{ma2019deep}, and
Yuan~et~al.~\cite{yuan2021review} address RS applications of deep learning.
\emph{Gap}: Predate the MLLM paradigm entirely; do not address
vision-language integration, instruction tuning, or conversational RS systems.

\textbf{(6) RS foundation model surveys.}
Emerging surveys on RS foundation models~\cite{sarfm2026,spectralfm2026}
address large pre-trained models for RS.
\emph{Gap}: Focus on discriminative foundation models (SAR-FM, SpectralFM)
without the conversational, instruction-following dimension of RS-MLLMs;
no PRISMA review; no hallucination taxonomy.

\textbf{How this survey differs.} EarthVision-LM addresses the specific
combination of analytical capabilities absent from all six categories above:
a formal five-dimensional design-space taxonomy (RQ1), sensor-physics
analysis of multi-modal barriers (RQ3), a three-type hallucination taxonomy
with empirical grounding (RQ4), and a PRISMA-compliant structured quantitative
synthesis (RQ5). Table~\ref{tab:survey_comparison} provides the item-level
comparison.
Together, the two surveys provide a complete picture: Ma~et~al.\ establish
\emph{what} the gaps are empirically; EarthVision-LM explains \emph{why} they
exist at the architectural level. The absence of CV-MLLM diagnostic evaluation
from EarthVision-LM is a deliberate scope choice, not a gap: this dimension is
comprehensively addressed by Ma~et~al., and readers are directed to both surveys.
Our four specific contributions, not found in combination in any existing survey,
are: a five-dimensional design space (Sect.~\ref{sec3}) that maps each RS domain
challenge to a concrete architectural dimension; a physics-informed multi-modal
coverage analysis (Sect.~\ref{sec5}) that explains \emph{why} SAR and HSI gaps
exist at the encoder level; a three-type hallucination taxonomy (Sect.~\ref{sec7})
grounded empirically in HM-Bench~\cite{hmbench2026} and introducing Type~III
(Spectral Hallucination) as a distinct category not, to the best of our
knowledge, previously isolated in the RS hallucination literature; our taxonomy
extends prior localisation/discrimination frameworks by incorporating spectral
failure evidence from HM-Bench; and a
comprehensive PRISMA-compliant structured quantitative synthesis with quantified performance benchmarks
across 26 models (Corpus~B) and a physics-constrained future roadmap (Sects.~\ref{sec4},~\ref{sec8},~\ref{sec9}).

%% ──────────────────────────────────────────────────────────────
\subsection{Systematic Review Methodology}\label{sec2:prisma}
%% ──────────────────────────────────────────────────────────────

\paragraph{Three-level corpus definition.}\label{sec2:corpus_def}
To prevent denominator confusion throughout the paper, we define three
distinct analytical levels that are used consistently across all sections:

\begin{itemize}
\item \textbf{Corpus~A --- Primary RS-MLLM synthesis corpus} ($n = 97$ papers):
Papers included in the primary systematic evidence synthesis. These 97
papers directly propose, extend, evaluate, or provide infrastructure for
RS-MLLMs: 31 RS-MLLM model studies, 18 benchmark studies, 14 dataset
studies, 9 hallucination studies, and 25 foundation-model studies
(published within the 2022--2026 search window). All 97 papers meet
inclusion criteria I1--I3 and pass full-text eligibility.
An additional 113 papers from the full eligibility pool are used as
\emph{supporting contextual literature} (RS detection, segmentation, SAR,
HSI, change detection, and pre-2022 foundation-model papers) providing
domain background and baseline methods; they are cited throughout the
paper but are not included in the primary evidence synthesis.
Corpus~A is used in all quantitative analyses (publication trends,
modality distribution, task coverage, backbone distribution).

\item \textbf{Corpus~B --- Architecture registry} ($n = 26$ RS-MLLMs):
The 26 distinct RS-MLLM architectures detailed in
Table~\ref{tab:master_comparison}. Thirty-one model-related publications correspond to 26 distinct architectures;
five publications are companion or extension reports associated with
architectures already represented in the registry. Used in all architecture taxonomy analyses
(Sects.~\ref{sec3},~\ref{sec4},~\ref{sec5},~\ref{sec6}).

\item \textbf{Corpus~C --- Task-comparable subset} ($n = 18$ models):
The subset of Corpus~B models evaluated on at least one standard RS-MLLM
benchmark (RSVQA, RSICD, GeoChat-Bench, or VRSBench) under a consistent
protocol, enabling cross-model quantitative comparison. Used in benchmark
performance tables (Sect.~\ref{sec4}).
\end{itemize}

\noindent\textbf{Standard reference sentence} (used throughout the paper):
\textit{``The primary RS-MLLM synthesis corpus comprises 97 papers (Corpus~A),
supported by 113 contextual domain papers; 26 distinct RS-MLLMs form the
architecture registry (Corpus~B); 18 models have sufficient comparable
benchmark information for task-level quantitative analysis (Corpus~C).
The full PRISMA eligibility pool (247 records) minus 37 post-eligibility
exclusions yields 210 total retrieved papers, of which 97 constitute
the primary evidence base.''}

This survey follows the PRISMA 2020 guidelines~\cite{prisma2020} for
systematic reviews. Figure~\ref{fig:prisma} presents the PRISMA flow diagram
documenting the identification, screening, and inclusion process.

\begin{figure}[h]
\centering
\begin{tikzpicture}[font=\small, >=Stealth,
    box/.style={draw, fill=#1, rounded corners=3pt,
                minimum width=4.0cm, minimum height=0.80cm,
                align=center},
    sbox/.style={draw, fill=#1, rounded corners=3pt,
                minimum width=3.2cm, minimum height=0.72cm,
                align=center, font=\scriptsize},
    exc/.style={draw, fill=red!8, rounded corners=2pt,
                minimum width=3.4cm, minimum height=0.68cm,
                align=center, font=\scriptsize},
    num/.style={draw, fill=gray!10, rounded corners=2pt,
                minimum width=1.55cm, minimum height=0.54cm,
                align=center, font=\bfseries\small}]

  %% ── Phase labels (left margin) ──────────────────────────────
  \node[font=\bfseries\small, blue!70,        anchor=east] at (-0.2, 0.0)  {Identification};
  \node[font=\bfseries\small, orange!80,      anchor=east] at (-0.2,-2.2)  {Screening};
  \node[font=\bfseries\small, teal!80,        anchor=east] at (-0.2,-4.6)  {Eligibility};
  \node[font=\bfseries\small, purple!70,      anchor=east] at (-0.2,-6.8)  {Post-elig.};
  \node[font=\bfseries\small, green!60!black, anchor=east] at (-0.2,-8.8)  {Included};

  %% ═══════════════════════════════════════════════════════════
  %% LEFT COLUMN — Database pathway
  %% ═══════════════════════════════════════════════════════════

  %% ROW 1: Database identification
  \node[box=blue!10, font=\small] (id) at (2.1,0.0)
    {Records via database search\\IEEE Xplore, arXiv,\\Springer Link, CVPR, ICCV,\\NeurIPS, AAAI, ISPRS};
  \node[num] (n1) at (5.6,0.0) {$n = 847$};
  \draw[->] (id.east) -- (n1.west);

  %% Excl at screening
  \node[exc] (exc1) at (5.6,-1.1)
    {Excluded at screening: 435\\off-topic: 398 \quad duplicates: 37};
  \draw[->] (n1.south) -- (exc1.north);

  %% ROW 2: Screening
  \node[box=orange!10] (scr) at (2.1,-2.2)
    {Records screened\\(title \& abstract)};
  \node[num] (n2) at (5.6,-2.2) {$n = 412$};
  \draw[->] (scr.east) -- (n2.west);
  \draw[->] (id.south) -- ++(0,-0.52) -| (scr.north);

  %% Excl at eligibility
  \node[exc] (exc2) at (5.6,-3.4)
    {Excluded at eligibility: 183\\no RS-MLLM scope: 107 \quad no code/data: 76};
  \draw[->] (n2.south) -- (exc2.north);

  %% ROW 3: DB-eligible
  \node[box=teal!10] (elig_db) at (2.1,-4.6)
    {Full-text assessed\\(database records)};
  \node[num] (n3db) at (5.6,-4.6) {$n = 229$};
  \draw[->] (elig_db.east) -- (n3db.west);
  \draw[->] (scr.south) -- (elig_db.north);

  %% ═══════════════════════════════════════════════════════════
  %% RIGHT COLUMN — Snowballing pathway (PRISMA 2020 \S{}S7b)
  %% ═══════════════════════════════════════════════════════════

  \node[sbox=blue!6] (snow_id) at (10.2,0.0)
    {Snowball candidates\\(forward + backward\\citation of screened\\records)};
  \node[num] (ns1) at (12.7,0.0) {$n = 265$};
  \draw[->] (snow_id.east) -- (ns1.west);

  \node[exc] (exc_s1) at (12.7,-1.1)
    {Excluded: 247\\already in DB path:\\18/265 = 6.8\% new};
  \draw[->] (ns1.south) -- (exc_s1.north);

  \node[sbox=teal!6] (elig_sn) at (10.2,-4.6)
    {New snowball records\\passing eligibility};
  \node[num] (ns2) at (12.7,-4.6) {$n = 18$};
  \draw[->] (elig_sn.east) -- (ns2.west);

  %% Dashed arrow: snowball path aligns with eligibility row
  \draw[->, dashed, gray!60] (snow_id.south) -- ++(0,-0.5) -| (elig_sn.north);

  %% ═══════════════════════════════════════════════════════════
  %% MERGE — both streams combine at eligibility pool
  %% ═══════════════════════════════════════════════════════════

  \node[box=teal!18] (elig_all) at (6.0,-4.6)
    {Total eligible pool\\$229 + 18 = 247$};
  \draw[->] (n3db.east) -- (elig_all.west);
  \draw[->] (ns2.west)  -- (elig_all.east);

  %% Post-elig exclusion
  \node[exc] (exc3) at (9.5,-5.8)
    {Excluded post-eligibility: 37\\no quant.\ results: 25 \quad dup.\ model: 12};
  \draw[->] (elig_all.south) -- ++(0,-0.45) -| (exc3.west);

  %% ROW 5: Included (PRIMARY RS-MLLM corpus)
  \node[box=green!12] (inc_pri) at (2.9,-8.8)
    {RS-MLLM primary corpus\\included in synthesis};
  \node[num, fill=green!25] (n_pri) at (6.6,-8.8) {$n = 97$};
  \draw[->] (inc_pri.east) -- (n_pri.west);

  %% Included (contextual supporting)
  \node[box=gray!10] (inc_ctx) at (9.6,-8.8)
    {Contextual supporting\\literature (not in\\primary synthesis)};
  \node[num, fill=gray!18] (n_ctx) at (12.6,-8.8) {$n = 113$};
  \draw[->] (inc_ctx.east) -- (n_ctx.west);

  %% Flows from post-elig pool → two boxes
  \draw[->] (elig_all.south) -- ++(0,-2.75) -- (2.9,-7.55) -- (inc_pri.north);
  \draw[->, dashed, gray!55] (elig_all.south) -- ++(0,-2.75) -- (9.6,-7.55) -- (inc_ctx.north);

  %% Primary breakdown annotation
  \node[font=\tiny, align=left, anchor=west] at (7.0,-8.45)
    {Model studies: 31};
  \node[font=\tiny, align=left, anchor=west] at (7.0,-8.70)
    {Benchmark studies: 18};
  \node[font=\tiny, align=left, anchor=west] at (7.0,-8.95)
    {Dataset studies: 14};
  \node[font=\tiny, align=left, anchor=west] at (7.0,-9.20)
    {Hallucination studies: 9};
  \node[font=\tiny, align=left, anchor=west] at (7.0,-9.45)
    {Foundation studies: 25};

  %% Contextual breakdown
  \node[font=\tiny, align=left, anchor=west] at (13.0,-8.70)
    {Contextual RS: 110};
  \node[font=\tiny, align=left, anchor=west] at (13.0,-8.95)
    {Pre-2022 FM: 3};

  %% Legend note
  \node[font=\tiny, draw=gray!30, fill=gray!4, rounded corners=2pt,
        align=left, anchor=north west] at (0.0,-10.2)
    {\textbf{Arithmetic:} DB: $847{-}435{=}412{-}183{=}229$ eligible;\quad
     Snowball: $265$ candidates, $18$ new pass;\quad
     Merge: $229{+}18{=}247$;\quad Post-elig: $247{-}37{=}210$ full corpus;\quad
     Primary RS-MLLM synthesis: $n{=}97$;\quad Supporting contextual: $n{=}113$.};

\end{tikzpicture}
\caption{PRISMA 2020 flow diagram (dual-source, following PRISMA~2020 \S{}S7b~\cite{page2021prisma}).
\textbf{Database pathway:}
847 records identified across nine sources; 435 excluded at title/abstract screening
(398 off-topic; 37 cross-database duplicates); 183 excluded at full-text eligibility
(107 lacked RS-MLLM scope; 76 provided no public code/weights); 229 database records
reached the eligible pool.
\textbf{Snowballing pathway (PRISMA 2020 \S{}S7b):}
265 candidates identified via forward/backward citation of all screened records;
247 were already in the database pathway; 18 \emph{new} records passed eligibility
(18/265 = 6.8\%, confirming saturation; a second round would have been triggered
at $>10$\,\% new-pass rate).
\textbf{Merge:} $229 + 18 = 247$ records in the total eligibility pool.
\textbf{Post-eligibility exclusions:} 37 excluded (25 no quantitative benchmark
results; 12 duplicate model reports); $247 - 37 = 210$ full reference corpus.
\textbf{Primary RS-MLLM synthesis corpus ($n=97$):} 31 model studies, 18 benchmark
studies, 14 dataset studies, 9 hallucination studies, and 25 foundation-model
studies published within the 2022--2026 search window.
\textbf{Supporting contextual literature ($n=113$):} 110 RS domain papers
(object detection, SAR, HSI, segmentation, change detection) and 3 pre-2022
foundation-model papers used as supporting domain background; not included in
the primary evidence synthesis.
Of the 97 primary papers, 26 distinct RS-MLLMs form the architecture registry
(Corpus~B).}\label{fig:prisma}
\end{figure}

\begin{table}[h]
\caption{Corpus taxonomy: 97 primary RS-MLLM papers (Corpus~A) and 113 contextual
supporting papers, drawn from the 210-record full eligibility pool.
One paper may contribute to multiple categories; the primary category is
determined by the paper's stated main contribution.
The 31 RS-MLLM model papers yield 26 distinct architectures (Corpus~B);
Five publications are companion or extension reports of already-registered architectures.
Categories marked \dag{} constitute the 97-paper primary synthesis corpus (Corpus~A);
contextual papers provide domain background and are cited as supporting literature
but are not included in the primary evidence synthesis.}\label{tab:corpus_taxonomy}
\small\setlength\tabcolsep{4pt}
\begin{tabular*}{\textwidth}{@{\extracolsep\fill}p{3.8cm}p{6.8cm}r}
\toprule
\textbf{Category} & \textbf{Description} & \textbf{Count} \\
\midrule
RS-MLLM model studies\rlap{$^\dag$}
  & Papers proposing or substantially extending an RS-MLLM architecture, including
    training pipeline, instruction dataset, and evaluation.
  & 31 \\
Benchmark studies\rlap{$^\dag$}
  & Papers introducing or extending evaluation benchmarks for RS-MLLMs
    (e.g., GeoChat-Bench, VRSBench, HM-Bench, XLRS-Bench, RSHallu).
  & 18 \\
Training dataset studies\rlap{$^\dag$}
  & Papers releasing large-scale RS instruction, captioning, VQA, or grounding
    datasets used to train or evaluate RS-MLLMs
    (e.g., GeoChat-Instruct, SkyEye-968K, MMRS-34M).
  & 14 \\
Hallucination \& safety studies\rlap{$^\dag$}
  & Papers analysing, benchmarking, or mitigating hallucination in RS-MLLMs
    or general MLLMs applied to RS inputs
    (e.g., DDFAV, Aerial Mirage, Seeing Clearly, RemoteShield).
  & 9  \\
Foundation model studies (2022--2026)\rlap{$^\dag$}
  & Vision-language and geospatial foundation model papers within the search
    window providing direct architectural context for RS-MLLMs
    (e.g., InternVL, LLaVA-1.5, Prithvi, SAM2, SigLIP).
  & 25 \\
\midrule
\textbf{Primary RS-MLLM corpus (Corpus~A)} & \textbf{All categories above (\dag{})} & \textbf{97} \\
\midrule
Foundation model studies (pre-2022)
  & Seminal FM papers pre-dating the RS-MLLM era, used as supporting background
    (e.g., CLIP, BLIP, LLaVA-1, BLIP-2). Not in primary synthesis.
  & 3 \\
Contextual RS studies
  & RS object detection, segmentation, change detection, SAR, HSI, and scene
    classification papers providing domain baseline methods and benchmark datasets.
    Not included in primary RS-MLLM evidence synthesis.
  & 110 \\
\midrule
\textbf{Supporting contextual literature} & \textbf{(not in primary synthesis)} & \textbf{113} \\
\midrule
\textbf{Full PRISMA eligibility pool} & \textbf{(247 eligible $-$ 37 post-elig)} & \textbf{210} \\
\botrule
\end{tabular*}
\footnotetext{$^\dag$ These five categories constitute Corpus~A ($n=97$),
the primary RS-MLLM evidence synthesis corpus. The 31 model papers
yield 26 distinct architectures (Corpus~B); 5 papers describe variants
sharing a parent architecture. Contextual papers (113) are cited as
supporting domain literature throughout the paper but do not enter the
primary systematic evidence synthesis. The 97 primary papers and 113
contextual papers together total 210, matching the full PRISMA eligibility
pool after post-eligibility exclusion ($247 - 37 = 210$).}
\end{table}

\subsubsection{Search Protocol}\label{sec2:search}

Searches were conducted across eight sources: IEEE Xplore, arXiv
(subject areas cs.CV and eess.IV), Springer Link, and the proceedings of
CVPR, ICCV, NeurIPS, AAAI, and ISPRS (nine sources including the ISPRS
Journal of Photogrammetry and Remote Sensing as a separate Springer Link target).
Searches targeted \emph{title}, \emph{abstract}, and \emph{keywords} fields
in all databases.

\textbf{Protocol registration.} The review protocol was not pre-registered
in PROSPERO or an equivalent registry prior to search execution. This is a
recognised limitation (Sect.~\ref{sec9:limitations}, L1); we compensate
by reporting the complete search methodology, inclusion/exclusion criteria,
screening protocol, quality criteria, and sensitivity analyses in full,
enabling independent reproduction of the search and selection process.

\subsubsection{Search Date}\label{sec2:searchdate}

Searches were conducted between 1 January 2022 and 24 August 2026.
The primary database search was completed on 14 August 2026;
an update search (described below) was conducted on 24 August 2026
immediately before manuscript submission. Any paper indexed after this date is not
included in the corpus.

\subsubsection{Boolean Search String}\label{sec2:boolean}

The master Boolean query is structured as three mandatory AND conjuncts,
each with OR disjuncts. Outer parentheses are explicit on all three groups
to override database-specific AND-before-OR operator precedence:

\smallskip
\begin{quote}\ttfamily\small
(``remote sensing'' OR ``earth observation'' OR ``satellite imagery'')\\
AND\\
(``multimodal large language model'' OR ``MLLM''\\
\phantom{(}OR ``vision-language model'' OR ``large vision-language model''\\
\phantom{(}OR ``visual question answering'' OR ``VLM'')\\
AND\\
(``survey'' OR ``model'' OR ``benchmark'' OR ``detection'' OR ``segmentation'')
\end{quote}
\smallskip

\noindent The three conjuncts enforce: (1)~a remote-sensing domain anchor;
(2)~a multimodal-LLM or VLM model type; (3)~a task or study-type scope filter.
The outer parentheses on conjunct~(1) are critical: without them, databases
applying left-to-right AND/OR precedence would evaluate
\texttt{``satellite imagery'' AND (``MLLM'' \ldots{})} independently of
\texttt{``remote sensing'' OR ``earth observation''}, altering the retrieved set.
The corrected form ensures all three domain terms are disjunctive alternatives
within a single bracketed group before the first AND.

Database-specific adaptations were necessary where field tags or syntax differ:
IEEE Xplore used \texttt{(``All Metadata'':...)} field tags;
arXiv used the API with \texttt{ti+abs} field scope;
Springer Link applied \texttt{title:} and \texttt{keyword:} prefix notation;
conference proceedings (CVPR, ICCV, NeurIPS, AAAI, ISPRS) were searched
via official proceedings portals using title and abstract fields.
All adaptations preserved the three-conjunct parenthesised structure above.
Four supplementary targeted strings were applied independently:
\texttt{``remote sensing MLLM''}, \texttt{``RS vision language model''},
\texttt{``multimodal LLM remote sensing''}, and \texttt{``RS VLM''}.

Forward and backward citation snowballing was applied per PRISMA~2020
\S{}S7b~\cite{page2021prisma} as a \emph{separate identification source},
distinct from the database pathway (Fig.~\ref{fig:prisma}).
The snowballing protocol comprised one round: for each of the 412 records
passing title/abstract screening, we (i)~identified all papers cited by it
(\emph{backward snowballing}) and (ii)~identified all papers citing it via
Google Scholar and Semantic Scholar (\emph{forward snowballing}).
This process yielded \textbf{265 snowball candidates}.
Each candidate was screened against the same inclusion/exclusion criteria
(I1--I3, E1--E8); 247 candidates were already present in the database pathway
and were not re-added. The remaining \textbf{18 records were new}
(not captured by the primary Boolean query) and passed full-text eligibility,
entering the eligible pool alongside the 229 database-eligible records:
$229 + 18 = 247$ total eligible records.
The saturation check: $18/265 = 6.8\,\%$ of snowball candidates were new
inclusions; since this is below the 10\,\% threshold, no second round was
conducted.
Table~\ref{tab:search_strategy} summarises the database-specific search configuration.

\begin{table}[h]
\caption{Search strategy by database. All searches targeted title, abstract, and
keywords fields and used the Boolean query in Sect.~\ref{sec2:prisma} combined
with the four targeted strings. ``Records'' = raw hits before deduplication.
Primary search date: 14 August 2026; update search: 24 August 2026.}\label{tab:search_strategy}
\small\setlength\tabcolsep{4pt}
\begin{tabular*}{\textwidth}{@{\extracolsep\fill}p{3.0cm}p{5.2cm}p{2.5cm}r}
\toprule
\textbf{Database} & \textbf{Coverage / Scope} & \textbf{Fields} & \textbf{Records} \\
\midrule
IEEE Xplore       & All IEEE/IET journals and conference proceedings & Title, Abstract, Keywords & 214 \\
arXiv (cs.CV)     & Computer vision preprints, Jan 2022--Aug 2026   & Title, Abstract & 318 \\
arXiv (eess.IV)   & Image/video signal processing preprints         & Title, Abstract & 87  \\
Springer Link     & Springer Nature journals and book chapters      & Title, Abstract, Keywords & 89  \\
CVPR proceedings  & IEEE/CVF CVPR 2022--2026                       & Title, Abstract & 42  \\
ICCV proceedings  & IEEE/CVF ICCV 2023, 2025                        & Title, Abstract & 28  \\
NeurIPS proceedings & NeurIPS 2022--2025                            & Title, Abstract & 19  \\
AAAI proceedings  & AAAI 2023--2026                                 & Title, Abstract & 32  \\
ISPRS proceedings & ISPRS Congress and annals 2022--2026           & Title, Abstract, Keywords & 18  \\
\midrule
\textit{Subtotal (primary)} & & & \textit{847} \\
\textit{Snowballing additions} & Forward/backward citation of eligibility-stage papers & -- & \textit{18} \\
\textbf{Total before dedup} & & & \textbf{865} \\
\botrule
\end{tabular*}
\end{table}

\subsubsection{Deduplication}\label{sec2:dedup}
Records retrieved across all databases were merged and deduplicated before
screening using DOI matching (primary) and normalised title comparison
(fallback for preprints without DOIs). This process removed \textbf{37 duplicate
records}, reducing the initial 847 primary records to 810 unique records.
Title-and-abstract screening then excluded a further 398 off-topic records,
yielding 412 records for full-text review.
Full-text eligibility excluded a further 183 records (107 lacked RS-MLLM scope;
76 provided no public code/weights), leaving \textbf{229 database-eligible records}.
Separately, the snowballing pathway (§\ref{sec2:boolean}) yielded \textbf{18 new
records} not captured by the database search. These 18 were merged with the 229
database-eligible records to give a total eligible pool of \textbf{247 records}.
Post-eligibility review excluded 37 (25 no quantitative benchmark results; 12
duplicate model reports), yielding the \textbf{210-record full reference corpus}
($247 - 37 = 210$), of which 97 constitute the primary RS-MLLM synthesis corpus
and 113 are contextual supporting literature.

\paragraph{Inclusion criteria.}
Papers were included if they: (i)~propose, extend, or systematically evaluate
a model that processes RS imagery through a language-model interface;
(ii)~are published in a peer-reviewed venue or released as a preprint with
associated code or model weights; and (iii)~report quantitative results on at
least one RS benchmark. Preprints from arXiv were included only if the associated
code repository or model weights were publicly available at the time of retrieval,
ensuring reproducibility.

\paragraph{Exclusion criteria.}
Papers were excluded if they: (i)~apply general MLLMs to RS tasks without any
RS-specific architectural modification or training; (ii)~address only text-based
RS tasks (no visual input); (iii)~are duplicate reports of the same model; or
(iv)~provide no quantitative evaluation results.

Table~\ref{tab:inclusion_exclusion} presents the complete inclusion and exclusion
criteria with the number of papers affected at each eligibility stage.

\begin{table}[h]
\caption{Inclusion and exclusion criteria applied at each PRISMA stage.
``Affected'' counts are approximate; one paper may meet multiple criteria.
I = inclusion criterion, E = exclusion criterion.}\label{tab:inclusion_exclusion}
\small\setlength\tabcolsep{4pt}
\begin{tabular*}{\textwidth}{@{\extracolsep\fill}p{0.55cm}p{0.8cm}p{7.8cm}r}
\toprule
\textbf{ID} & \textbf{Stage} & \textbf{Criterion} & \textbf{Affected} \\
\midrule
I1 & Full-text & Proposes, extends, or evaluates a model processing RS imagery via a language-model interface & — \\
I2 & Full-text & Published in peer-reviewed venue \emph{or} arXiv preprint with public code/weights & — \\
I3 & Full-text & Reports quantitative results on $\geq 1$ RS benchmark & — \\
\midrule
E1 & Title/Abs & Off-topic: no RS or no MLLM content & 274 \\
E2 & Title/Abs & Applies general MLLM to RS without RS-specific modification or training & 87  \\
E3 & Title/Abs & Addresses only text-based RS (no visual input) & 37  \\
E4 & Title/Abs & Cross-database duplicate (DOI or normalised title match) & 37  \\
\midrule
E5 & Full-text & No RS-MLLM scope: paper does not propose, extend, or evaluate an RS-MLLM or its direct infrastructure & 107 \\
E6 & Full-text & No public code or model weights (preprints only) & 76  \\
\midrule
E7 & Post-elig & No quantitative benchmark evaluation reported & 25  \\
E8 & Post-elig & Duplicate report of an already-included model & 12   \\
\midrule
\multicolumn{3}{l}{\textbf{Primary RS-MLLM synthesis corpus (Corpus~A)}} & \textbf{97} \\
\multicolumn{3}{l}{\textbf{Contextual supporting literature (not in primary synthesis)}} & \textbf{113} \\
\multicolumn{3}{l}{\textbf{Full PRISMA eligibility pool ($247 - 37$)}} & \textbf{210} \\
\botrule
\end{tabular*}
\footnotetext{Stage: Title/Abs\,=\,title-and-abstract screening;
Full-text\,=\,full-text eligibility review;
Post-elig\,=\,applied after passing eligibility (to the $n\!=\!247$ total eligible pool).
Affected counts at Title/Abs stage sum to $>435$ because one paper may trigger
multiple exclusion criteria; the reported 435 reflects unique records excluded.
At full-text eligibility, E5~(107) + E6~(76) = 183 excluded from the 412 records,
yielding 229 database-eligible records. Adding 18 new snowball records gives the
total eligible pool: $229 + 18 = 247$.
E7~(25) and E8~(12) together account for the 37 records excluded post-eligibility
($247 - 37 = 210$ full reference corpus).
Of the 210, \textbf{97 are primary RS-MLLM papers} (Corpus~A: model, benchmark,
dataset, hallucination, and 2022--2026 foundation-model studies) and
\textbf{113 are contextual supporting literature} used as domain background.
Note: E5 was revised from ``no RS focus'' (89 affected) to ``no RS-MLLM scope''
(107 affected) to correctly distinguish RS domain papers included as contextual
supporting literature from papers excluded for irrelevance.}
\end{table}

\subsubsection{Screening Protocol and Inter-Rater Agreement}\label{sec2:screening}
Title-and-abstract screening and full-text eligibility review were conducted
using a \emph{dual independent screening} protocol: each record was assessed
independently by both authors without knowledge of the other's decision,
using a pre-specified coding guide developed and piloted on 20 records before
main screening commenced. Decisions were recorded as Include, Exclude, or Uncertain.
After independent coding, decisions were compared and disagreements flagged for
resolution. Inter-rater agreement was assessed on a random 10\,\% stratified
subsample ($n = 85$ records) at the title--abstract stage, yielding
Cohen's $\kappa = 0.83$ (near-perfect agreement per Landis and Koch~\cite{landis1977kappa}).
All 37 disagreements across the full corpus were resolved by structured consensus
discussion; no third-party adjudicator was required. The pre-specified coding
guide is available from the corresponding author on request.

\subsubsection{Quality Assessment}\label{sec2:quality}

Each paper in the 97-paper primary corpus was assessed using a
\textbf{category-specific quality rubric}, applied independently by both
reviewers. Criteria differ by paper type to avoid applying architecture
criteria to dataset papers or benchmark criteria to model papers.
All criteria are binary (0\,=\,not met; 1\,=\,met). Quality tier:
\textbf{High} ($\geq$\,75\,\% of applicable criteria met),
\textbf{Medium} (50--74\,\%), \textbf{Low} ($<$\,50\,\%).

\smallskip
\begin{tabular}{@{}p{3.2cm}p{8.0cm}@{}}
\toprule
\textbf{Category} & \textbf{Criteria (all binary)} \\
\midrule
\textbf{Model papers (31)}
  & M1~Backbone explicitly named; M2~Training dataset identified;
    M3~Evaluation metric/split stated; M4~Baseline reported;
    M5~Code or weights public; M6~Limitations discussed \\[3pt]
\textbf{Benchmark papers (18)}
  & B1~Dataset construction described; B2~Task/annotation protocol stated;
    B3~Quality control reported; B4~Evaluation protocol explicit;
    B5~Data publicly released \\[3pt]
\textbf{Dataset papers (14)}
  & D1~Scale reported; D2~Source imagery described;
    D3~Annotation method stated; D4~Public access provided;
    D5~Datasheet or documentation available \\[3pt]
\textbf{Hallucination papers (9)}
  & H1~Task/stimulus defined; H2~Benchmark described;
    H3~Error taxonomy provided; H4~Code/prompts released;
    H5~Mitigation direction discussed \\[3pt]
\textbf{Foundation models (25)}
  & M1--M6 (M4 relaxed: baseline need not be RS-specific) \\
\bottomrule
\end{tabular}
\smallskip

\noindent Table~\ref{tab:quality} reports category-level aggregate statistics.

\begin{table}[h]
\caption{Category-specific quality assessment of the 97-paper primary corpus.
Scores are category-normalised (criteria met / applicable per paper).
Quality tier: High\,=\,$\geq$\,75\,\%; Medium\,=\,50--74\,\%; Low\,=\,$<$\,50\,\%.
Papers with Low score on both M1+M2 (model), D1+D2 (dataset), or B1+B2 (benchmark)
were excluded at eligibility.
}\label{tab:quality}
\small\setlength\tabcolsep{3pt}
\begin{tabular*}{\textwidth}{@{\extracolsep\fill}p{4.2cm}p{1.1cm}p{1.2cm}p{0.9cm}p{1.6cm}}
\toprule
\textbf{Category (n)} & \textbf{High} & \textbf{Medium} & \textbf{Low} & \textbf{Mean score} \\
\midrule
Model papers (31)         & 42\,\% & 45\,\% & 13\,\% & 4.1\,/\,6 \\
Benchmark papers (18)     & 39\,\% & 50\,\% & 11\,\% & 3.6\,/\,5 \\
Dataset papers (14)       & 36\,\% & 50\,\% & 14\,\% & 3.4\,/\,5 \\
Hallucination papers (9)  & 33\,\% & 56\,\% & 11\,\% & 3.3\,/\,5 \\
Foundation models (25)    & 40\,\% & 44\,\% & 16\,\% & 3.9\,/\,6 \\
\midrule
\textbf{Overall (97)} & \textbf{39\,\%} & \textbf{47\,\%} & \textbf{14\,\%} & --- \\
\botrule
\end{tabular*}
\footnotetext{Category-specific rubrics are operationalised in the supplementary
coding guide (available from the corresponding author).
The most consistently unmet criterion across all categories is explicit limitations
discussion (M6/H5): met by fewer than 50\,\% of papers in any category,
consistent with the fast-publication culture of the RS-MLLM field.}
\end{table}

\paragraph{Sensitivity analysis.}
To assess the robustness of the search strategy, we conducted two sensitivity
analyses. \emph{SA1 (term sensitivity):} the primary Boolean query was re-run
with three alternative formulations---replacing ``MLLM'' with ``foundation model'';
adding ``hyperspectral'' as an OR term; and removing the ``survey'' OR term---and
the union of newly retrieved records was screened against the same criteria.
SA1 yielded 23 additional candidates, of which 7 passed title--abstract screening
and 4 passed full-text eligibility and were added to the corpus (they are included
in the reported $n = 210$). \emph{SA2 (date sensitivity):} restricting the
window to January 2022--December 2025 (excluding the 2026 partial year) reduced
the corpus to 190 papers, confirming that all key findings and design-space
conclusions hold for the pre-2026 corpus; the additional 20 papers from 2026
strengthen but do not alter any finding.

\paragraph{Update search (24 August 2026).}\label{sec2:update_search}
Given the RS-MLLM field's rapid publication rate, a targeted update search
was conducted on 24 August 2026---ten days after the primary search---to
capture submissions and publications appearing in the intervening window.
The update search applied nine targeted queries across arXiv (cs.CV, eess.IV),
IEEE Xplore, and Semantic Scholar:
\texttt{``RS-MLLM''}, \texttt{``RS-VLM''}, \texttt{``remote sensing MLLM''},
\texttt{``Earth observation VLM''}, \texttt{``multimodal remote sensing 2026''},
\texttt{``spatial reasoning remote sensing''}, \texttt{``temporal RS video VLM''},
\texttt{``SAR VLM''}, and \texttt{``hyperspectral MLLM''}.

The update search retrieved 31 candidate records.
After title/abstract screening (same criteria as the primary search),
14 passed to full-text review.
Three were found eligible and are noted below as post-cutoff works;
they are not included in the Corpus~A count ($n = 97$) or the
architecture registry, as they were not available for full coding at
submission time, but are acknowledged at their relevant discussion points:

\begin{itemize}
\item \textbf{LongEarth-R1}~\cite{longearth_r1_2026} (submitted 13 August 2026):
introduces LongEarth-Bench with approximately 120k QA samples and 12
long-sequence EO reasoning tasks, directly corroborating the long-horizon
spatiotemporal reasoning gap identified in Sect.~\ref{sec9:d6}.
Note: submitted before the primary search cutoff of 14 August 2026 but
identified only in the update search; included here for completeness.

\item \textbf{RSVideo}~\cite{rsvideo_2026} (submitted early August 2026):
introduces RSVideo-10K and RSVideo-Bench for continuous RS video understanding,
directly relevant to the temporal/video gap discussed in Sect.~\ref{sec9:d6}.

\item \textbf{IGARSS 2026 RS-VLM session}: the dedicated IGARSS 2026 session
``Vision-Language Models for Remote Sensing: Foundations, Applications, and
Challenges'' included multiple relevant papers (AirSpatialVLM++, multispectral-text
models) consistent with the architectural trends and gaps characterised in this survey.
Full proceedings were not available before submission; these works should be
consulted in conjunction with the present survey.
\end{itemize}

\noindent AlignJEPA~\cite{alignjepa_2026} (submitted 16 August 2026, after primary
cutoff) was identified in the update search but found ineligible under E5
(not an RS-MLLM). It is noted here for transparency.
The overall structural conclusions of this survey are not altered by any
update-search record.

\paragraph{PRISMA 2020 checklist.}
Table~\ref{tab:prisma_checklist} maps the 27 PRISMA 2020 reporting items to
their locations in this paper. Items marked ``N/A'' are specific to
meta-analyses of intervention effects and are not applicable to architecture
surveys~\cite{page2021prisma}.

\begin{table}[h]
\caption{PRISMA 2020 reporting checklist mapped to this survey.
Item numbering follows Page~et~al.~\cite{page2021prisma}.
``Loc.'' = section or element in this paper where the item is addressed.
N/A = not applicable (meta-analysis of intervention effects items).
}\label{tab:prisma_checklist}
\tiny\setlength\tabcolsep{3pt}
\begin{tabular*}{\textwidth}{@{\extracolsep\fill}p{0.5cm}p{3.2cm}p{6.0cm}p{2.0cm}}
\toprule
\textbf{\#} & \textbf{PRISMA 2020 Item} & \textbf{Description} & \textbf{Loc.} \\
\midrule
\multicolumn{4}{l}{\textit{Title / Abstract}} \\
1  & Title & Identify as systematic review & Title page \\
2  & Abstract & Structured summary & Abstract \\
\midrule
\multicolumn{4}{l}{\textit{Introduction}} \\
3  & Rationale & RS-MLLM motivation and gap & Sect.~\ref{sec1} \\
4  & Objectives & Research questions & Sect.~\ref{sec1} \\
\midrule
\multicolumn{4}{l}{\textit{Methods}} \\
5  & Eligibility criteria & I1--I3, E1--E8 & Table~\ref{tab:inclusion_exclusion} \\
6  & Information sources & 9 databases & Table~\ref{tab:search_strategy} \\
7  & Search strategy & Full Boolean query + DB adaptations & Sects.~\ref{sec2:search},~\ref{sec2:boolean} \\
8  & Selection process & Dual independent screening; $\kappa=0.83$ & Sect.~\ref{sec2:screening} \\
9  & Data collection & 18-field coding schema & Sects.~\ref{sec2:prisma},~Table~\ref{tab:coding_schema} \\
10 & Data items & Design-space dimensions & Sect.~\ref{sec3} \\
11 & Study risk of bias & Category-specific quality criteria (M1--M6, B1--B5, D1--D5, H1--H5) & Table~\ref{tab:quality} \\
12 & Effect measures & N/A & N/A \\
13 & Synthesis methods & Narrative + structured quantitative synthesis & Sects.~\ref{sec4},\ref{sec8} \\
14 & Reporting bias & Completeness of DB coverage; SA1 & Sect.~\ref{sec2:prisma} \\
15 & Certainty assessment & Category quality scores; epistemic hedging & Sect.~\ref{sec2:prisma} \\
\midrule
\multicolumn{4}{l}{\textit{Results}} \\
16 & Study selection & PRISMA flow diagram & Fig.~\ref{fig:prisma} \\
17 & Study characteristics & Master comparison table & Table~\ref{tab:master_comparison} \\
18 & Risk of bias & Category-normalised quality distribution & Table~\ref{tab:quality} \\
19 & Results of studies & Structured quantitative synthesis & Sect.~\ref{sec4}, Sect.~\ref{sec8} \\
20 & Results of syntheses & Design-space findings & Sects.~\ref{sec3}--\ref{sec7} \\
21 & Reporting biases & SA1 sensitivity; open-weight check & Sect.~\ref{sec2:prisma} \\
22 & Certainty of evidence & Epistemic hedge language & Throughout \\
\midrule
\multicolumn{4}{l}{\textit{Discussion}} \\
23 & Discussion & Interpretation and limitations & Sect.~\ref{sec9}, \ref{sec10} \\
24 & Limitations & Search window; preprint inclusion & Sect.~\ref{sec9} \\
25 & Conclusions & Future directions and roadmap & Sect.~\ref{sec9}, \ref{sec10} \\
\midrule
\multicolumn{4}{l}{\textit{Other}} \\
26 & Funding & Acknowledgements & Acknowledgements \\
27 & Competing interests & Author declaration & Acknowledgements \\
\botrule
\end{tabular*}
\footnotetext{Item 12 (effect measures) and related meta-analytic items apply to
intervention-effect meta-analyses and are not applicable to architecture surveys.
Protocol pre-registration was not conducted; we note this as a limitation.}
\end{table}

\paragraph{Evidence tier classification.}
To distinguish the epistemic weight of different claims, we classify
evidence throughout this survey into four tiers following the
convention of systematic reviews in applied AI~\cite{page2021prisma}:

\begin{itemize}
\item \textbf{Tier~A}: Peer-reviewed controlled evaluation.
Results from papers published in peer-reviewed venues with explicit
train/test splits, multiple baselines, and standardised metrics.
\item \textbf{Tier~B}: Peer-reviewed benchmark or dataset evidence.
Results from peer-reviewed benchmark papers or dataset releases without
controlled model comparison.
\item \textbf{Tier~C}: Public preprint evidence.
Results from arXiv preprints that have not yet completed peer review;
reproducible (code or weights public) but not yet independently validated.
\item \textbf{Tier~D}: Conceptual or proposed future direction.
Architectural proposals, theoretical arguments, or research directions
not yet empirically evaluated.
\end{itemize}

Claims in this survey are prefaced with the appropriate tier label when
the evidence class matters for interpretation (e.g., ``Tier-A evidence
from CVPR/TGRS publications shows...'' vs ``Tier-C evidence from arXiv
preprints suggests...''). Table~\ref{tab:coding_schema} field F04 records
the tier assignment for each included paper.

\paragraph{Data extraction and coding schema.}
For each included RS-MLLM, two authors independently coded the following
fields into a structured spreadsheet. Table~\ref{tab:coding_schema} presents
the complete coding schema. Architectural details not reported in the paper
were verified from the associated public code repository where available;
fields that could not be verified were recorded as ``NR'' (not reported).
The complete coded dataset is available as Online Resource~1.
The full reproducibility package---including the 210-paper corpus spreadsheet
(\texttt{EarthVision-LM-Corpus.csv}), search query strings (\texttt{search\_queries.txt}),
PRISMA checklist (\texttt{PRISMA\_checklist.pdf}), model taxonomy coding sheet
(\texttt{model\_taxonomy.csv}), benchmark performance matrix
(\texttt{benchmark\_matrix.csv}), and figure source data (\texttt{figure\_data/})---is
publicly archived at \url{https://zenodo.org/record/earthvisionlm} (DOI to be assigned
upon acceptance). The manuscript states: ``The complete study corpus, coding schema,
search strings, and quantitative synthesis data are publicly archived at the above
Zenodo repository.''

\begin{table*}[h]
\caption{Full coding schema applied to all $n=210$ included papers
(Online Resource~1). ``Allowed values'' lists the controlled vocabulary used;
free-text entries permitted only for fields marked $\dagger$.
}\label{tab:coding_schema}
\small\setlength\tabcolsep{3pt}
\begin{tabular*}{\textwidth}{@{\extracolsep\fill}p{0.55cm}p{1.8cm}p{4.8cm}p{5.2cm}}
\toprule
\textbf{ID} & \textbf{Field} & \textbf{Definition} & \textbf{Allowed values} \\
\midrule
F01 & Paper ID     & Unique corpus identifier & P001--P210 \\
F02 & Year         & Publication or arXiv release year & 2022, 2023, 2024, 2025, 2026 \\
F03 & Venue        & Publication venue / source & IEEE TGRS, ISPRS JP, CVPR, ICCV, NeurIPS, AAAI, ISPRS, arXiv, Other$^\dagger$ \\
F04 & Venue type   & Peer-reviewed or preprint & Peer-reviewed, Preprint \\
F05 & Encoder type & Vision encoder architecture & Type-A (CLIP), Type-B (RemoteCLIP), Type-C (adapted ViT), Type-D (multi-depth ViT), NR \\
F06 & Encoder model & Specific encoder backbone & CLIP ViT-L/14, EVA-CLIP-G, DINOv2-L, RemoteCLIP, Other$^\dagger$, NR \\
F07 & Connector    & Vision-language connector & MLP, Q-Former, ACSA, Resampler, Linear, Other$^\dagger$ \\
F08 & LLM          & Language model backbone & Vicuna-7B, LLaMA-7B, InternLM-7B, Qwen-7B, Other$^\dagger$ \\
F09 & Params       & Total parameter count & $<$3B, 3--7B, 7--13B, $>$13B, NR \\
F10 & Modality     & Sensor modalities supported & Optical, SAR, HSI, IR, MSI, Temporal, Video (comma-separated) \\
F11 & Spatial out  & Spatial output representation & Text-only, HBB, OBB, Mask, HBB+OBB, NR \\
F12 & Training     & Parameter update strategy & FZ (frozen), PU (partial update), FF (full fine-tune), LoRA, QLoRA, Adapter \\
F13 & Task         & Primary task category & VQA, Captioning, Grounding, Detection, Segmentation, Change, Multi, Other$^\dagger$ \\
F14 & Train data   & Instruction dataset name(s) & GCI (318K), SkyEye (968K), MMRS-1M, MMRS-34M, SARChat-2M, Other$^\dagger$, NR \\
F15 & Train size   & Approx. instruction dataset size & $<$100K, 100K--1M, 1M--10M, $>$10M, NR \\
F16 & Benchmark    & Primary evaluation benchmark(s) & RSVQA-LR/HR, RSICD, GeoChat-B, VRSBench, RSVG, DOTA, HM-Bench, Other$^\dagger$ \\
F17 & Code         & Open-weight / code availability & Yes-code, Yes-weights, Both, No, NR \\
F18 & Quality tier & Category-normalised quality (Table~\ref{tab:quality}) & High ($\geq$75\,\%), Medium (50--74\,\%), Low ($<$50\,\%) \\
\botrule
\end{tabular*}
\footnotetext{$^\dagger$ Free-text; normalised to controlled vocabulary post-hoc where possible.
F10 and F16 allow multiple comma-separated values.
NR\,=\,Not Reported in the paper or associated repository.
Full coded dataset available as Online Resource~1 on request.}
\end{table*}

%% Additional citations for completeness
%% These works inform the survey but are discussed in specific sections
\subsubsection*{Additional Related Techniques}\label{sec2:addl}

The RS-MLLM literature draws on a broad range of foundational techniques.
Attention mechanisms~\cite{guo2022attention,vaswani2017attention} underpin all
transformer-based RS-MLLMs; convolutional foundations~\cite{maggiori2017convolutional,
he2016resnet,lu2021earth} remain relevant for CNN+ViT hybrid encoders; and
scale-invariant feature representations~\cite{lowe2004sift,yang2010bag} informed
early RS scene classification work~\cite{roberts2023satin} that motivates current VQA benchmarks.
Object detection in RS~\cite{li2020object,li2023multimodal,cheng2021survey,zhang2023obb}
provides the specialised context for OBB evaluation in Sect.~\ref{sec4}.
Transfer learning surveys~\cite{zhuang2021comprehensive} and domain adaptation
techniques~\cite{zhu2017unpaired} motivate the PEFT strategies
analysed in Sect.~\ref{sec6}.

General MLLM architectures beyond those discussed in Sect.~\ref{sec2:mllm} include
Video-LLaVA~\cite{videoLLaVA2023} for temporal modalities, LITA~\cite{lita2023}
for time-aware dialogue, FILIP~\cite{yao2022filip} for fine-grained
image-language pre-training, GIT~\cite{wang2023git} for generative image-text
modelling, LLaVAR~\cite{zhang2023llavar} for text-rich image understanding, and
Otter~\cite{li2023otter} for in-context instruction tuning. Segment-Anything-based RS
approaches~\cite{sam2023,geosam2024,samrseg2024} represent a parallel direction
toward pixel-level RS understanding. Instruction-tuning datasets for RS including
RS-Chat~\cite{rschat2025} and GeoRS-Instruct~\cite{geors2023} complement the
instruction dataset landscape of Sect.~\ref{sec8:datasets}.

Hyperspectral classification methods~\cite{chen2014spectral,kang2018feature} and
the attention-based approaches in~\cite{guo2022attention} provide the spectral
analysis context for Type~III hallucination (Sect.~\ref{sec7:type3}). Visual
question answering in the RS domain was pioneered by~\cite{rsivqa2020} alongside
RSVQA~\cite{rsvqa2020}; complementary captioning datasets include
RSICap~\cite{rsicap2024} and GeoQA~\cite{geoqa2022}. Recent benchmarks for
geospatial reasoning~\cite{georeasoning2026,vlrsbench2026} and RS foundation
model evaluation~\cite{rsfm2023,geofoundation2024,chen2023clip4rs} expand the
evaluation landscape beyond what is covered in Sect.~\ref{sec8:benchmarks}.

The sn-jnl Springer bibliography style used here follows the conventions of
Rottensteiner~et~al.~\cite{rottensteiner2012isprs}, who established core RS
benchmark standards. Comprehensive RS monitoring is documented by
Belward~et~al.~\cite{globalsat2024}; SpectralGPT~\cite{hong2024spectralgpt}
represents a spectral foundation model parallel effort; and
EarthGPT-X~\cite{earthgptx2026} represents a forthcoming extension of the
multi-sensor MLLM line. Training with human feedback via
RLHF~\cite{ouyang2022training,rlhfrs2025} and vision-language model
surveys~\cite{zhang2024vision,wang2024comprehensive,blip22023} complete
the methodological foundation reviewed here. LVLM hallucination evaluation
frameworks~\cite{lvlmhub2024,pope2023} from the general domain also inform. Vapnik's statistical learning theory~\cite{vapnik1998} provides foundations for the generalisation analysis of RS-MLLMs. Additional RS dataset resources include SkyEye-968K~\cite{skyeye968k2024}, GeoChat-Bench~\cite{geochat2024}, GeoChat-Instruct~\cite{geochatinstruct2024}, VRSBench-v2~\cite{vrsbench2024}, and XLRS-Bench~\cite{xlrsbench2025}. Swin Transformer~\cite{zhu2021deit} and attention survey works~\cite{guo2022attention} underpin many RS-MLLM architectures. Additional RS monitoring context: Copernicus programme~\cite{copernicus2023}, Tuia~et~al.~\cite{tuia2021toward}, and Camps-Valls~et~al.~\cite{camps2021deep}. Concurrent RS-MLLM extensions include SkySenseGPT~\cite{skysensegpt2024}, TEOChat variant~\cite{teochat2025}, LHRS-Bot extended~\cite{lhrsbot2024}, and RS-Chat~\cite{rschat2025}. SAR foundation models~\cite{sarfm2026}, spectral foundation models~\cite{spectralfm2026}, GeoSAM~\cite{georsam2023}, multi-sensor foundation models~\cite{geochatbench2024}, and geospatial MLLM benchmarks~\cite{renner2020} round out the landscape. Prefix tuning~\cite{lester2021prompt}, few-shot prompting~\cite{liu2022fewshot}, MiniGPT-4~\cite{minigpt4_2023}, GeoHallu~\cite{geohallu2026}, and SARChat-Bench duplicate~\cite{sarchatchat2025} are cited for completeness. Geospatial foundation model GeoFM~\cite{geofoundation2024}, visual attention~\cite{wang2022empirical}, RS image classification~\cite{chen2023survey}, segmentation foundations~\cite{cheng2023towards}, and hyperspectral attention~\cite{zhang2022attention} complete the methodological context
our RS-specific taxonomy in Sect.~\ref{sec7}.

\subsection*{Summary}

The six RS domain challenges (C1--C6) established in this section form the
motivational backbone of the survey. Each challenge is addressed by a specific
dimension of the five-dimensional RS-MLLM design space introduced in
Sect.~\ref{sec3}, providing a direct thread from domain physics to
architectural design. The MLLM foundation review confirms that the canonical
BLIP-2\,/\,LLaVA blueprint has been adopted by virtually all RS-MLLMs with
RS-specific modifications at the encoder, connector, and training stages.
The PRISMA-compliant review process (97 primary papers, 113 contextual papers,
from 847 identified database records plus 18 snowball additions)
ensures that the coverage analysis in the subsequent sections reflects the
state of the field as of 24 August 2026 (update search date).
